\documentclass[UTF8,12pt,a4paper]{ctexart}

\usepackage{amsmath}
\usepackage{cases}
\usepackage{cite}
\usepackage{graphicx}
\usepackage{enumerate}
\usepackage{algorithm}
\usepackage{caption} %\caption*需要
\usepackage[noend]{algpseudocode} %algorithmicx
\makeatletter
\renewcommand{\fnum@algorithm}{\fname@algorithm}
\makeatother
\renewcommand{\algorithmicrequire}{\textbf{Input:}}
\renewcommand{\algorithmicensure}{\textbf{Output:}}
\usepackage[margin=1in]{geometry}
\geometry{a4paper}
\usepackage{fancyhdr}
\pagestyle{fancy}
\fancyhf{}
\usepackage{xcolor}
\usepackage{listings}
\lstset{
	basicstyle=\ttfamily,                                % 传统上，代码打字机字体好些
	breaklines,
	tabsize=4,
	extendedchars=false,
	columns=fixed,
	%numbers=left,                                        % 在左侧显示行号
	numbers=none,
	frame=none,                                          % 不显示背景边框
	backgroundcolor=\color[RGB]{245,245,244},            % 设定背景颜色
	keywordstyle=\color[RGB]{40,40,255},                 % 设定关键字颜色
	numberstyle=\footnotesize\color{darkgray},           % 设定行号格式
	commentstyle=\it\color[RGB]{0,96,96},                % 设置代码注释的格式
	stringstyle=\rmfamily\slshape\color[RGB]{128,0,0},   % 设置字符串格式
	showstringspaces=false,                              % 不显示字符串中的空格
	xleftmargin=1em,
	%xrightmargin=1em,
	%aboveskip=-0.5em
	language=c++,                                        % 设置语言
}

% title依据是作业还是实验报告，以及是第几次作业/实验报告，
% 比如第1次分治算法作业，则命名为"算法原理 HW1 分治算法";
% 第1次分治算法实验报告，则命名为"算法原理 Lab1 分治算法"
\title{算法原理 HW3 作业}
\author{
	姓名：\underline{冯峻}~~~~~~
	学号：\underline{523031910148}~~~~~~}
\date{\today}
\pagenumbering{arabic}

\begin{document}

% 如下两行请按实际情况修改
\fancyhead[L]{冯峻}
\fancyhead[C]{贪心算法}
\fancyfoot[C]{\thepage}

\maketitle
%\tableofcontents
%\newpage


\begin{enumerate}[1.]
	\item 你作为某公司的经理，今天有$n$件事情$a_1,a_2,··· ,a_n$ 需要处理。对于每件事情$a_i$，其处理时长为$l_i$，即需要$l_i$时间进行处理；同时每件事情$a_i$具有截止时间$e_i$，即事情$a_i$必须在时间$e_i$之前得到处理。已知只要按照正确的顺序进行处理，今天的n件事情都可以在其截止时间之前得到妥善处理。注意，你每次只能处理一件事情，且处理过程中不能中断去处理别的事情。请证明，只要每次选择截止时间最早的事情进行处理，即可让今天的n件事情都可以在截止时间之前得到处理。

	\textbf{解：}

	我们只需要构造这样一种任务序列： $(a_{j1}, a_{j2}, \dots, a_{jn})$，满足$e_{j1}\leq e_{j2}\leq \dots e_{jn}$，然后证明该任务序列能够保证所有任务都在截止时间处理完即可。

	设存在一种可行任务序列$\sigma=(a_{i_1},a_{i_2},\dots,a_{i_n})$，\textbf{其中每个任务均能在其截止时间之前完成。}
	
	若$\sigma$不是按截止时间非递减的顺序排列，则必存在相邻两任务$a_{i_k},a_{i_{k+1}}$使得$e_{i_k}>e_{i_{k+1}}$。考虑交换它们的执行顺序，得到新任务序列$\sigma'=(a_{i_1},\dots,a_{i_{k+1}},a_{i_k},\dots,a_{i_n})$。

	设$a_{i_k}$与$a_{i_{k+1}}$在原任务序列中开始时刻为$t$，则原任务序列满足：
	\[
	\begin{cases}
	t+l_{i_k} \le e_{i_k}, \dots (1)\\
	t+l_{i_k}+l_{i_{k+1}} \le e_{i_{k+1}}. \dots (2)
	\end{cases}
	\]
	交换后，$a_{i_{k+1}}$完成于$t+l_{i_{k+1}}$，$a_{i_k}$完成于$t+l_{i_{k+1}}+l_{i_k}$。我么们需要证明：
	\[
	\begin{cases}
	t+l_{i_{k+1}} \le e_{i_{k+1}}, \dots (3)\\
	t+l_{i_{k+1}}+l_{i_{k}} \le e_{i_{k}}. \dots (4)
	\end{cases}
	\]
	由(2)式易知：$t+l_{i_{k+1}} < t+l_{i_k}+l_{i_{k+1}} \le e_{i_{k+1}}$

	由于$e_{i_{k+1}}<e_{i_k}$，显然有：$t+l_{i_{k+1}}+l_{i_k} \le e_{i_{k+1}} \le e_{i_k}$
	
	故交换后两任务仍能在截止时间前完成，新任务序列仍然可行。

	重复此交换操作，直至所有相邻任务满足$e_{i_1}\le e_{i_2}\le\dots\le e_{i_n}$，即可得到一个按截止时间非递减排序的可行任务序列。

	\item 假设有 $n$ 个文件需要存储到 U 盘上，文件 $F_i$ 的尺寸为 $S_i$，已知 U 盘的容量为 $D$，且$D<\sum_{i=1}^{n}S_i$，即 U 盘不能放下所有的$n$个文件。

	1) 如果对文件按照其尺寸的升序排列，然后从最小文件开始，依次选择文件放入U盘，是否可使得U盘中存放最多数目的文件？如是，请证明；如不是，请给出反例。

	2) 如果对文件按照其尺寸的降序排列，然后从最大文件开始，依次选择文件放入U盘，是否可使得U盘空间的利用率最高，即剩余最少的空间？如是，请证明；如不是，请给出反例

	\textbf{1) 解：}
	是的，证明如下：
	我们只需要证明算法的贪心性质，即只需要证明每次都选择尺寸最小的文件，就能够得多文件个数最多的集合$P$。在下面的证明过程中，我们不妨设：
	$$S_1 \leq S_2 \leq S_3 \leq \dots \leq S_\text{n}$$
	用$|P|$表示集合$P$中\textbf{文件尺寸之和}
	\begin{enumerate}
		\item \textbf{1. 贪心选择性：必然存在某个优化解$P$包括文件$F_1$}
		
		我们假设集合$P$是一个优化解，并且$F_1 \notin P$。
		
		因为$P$是一个优化解，所以必然满足：$|P| \leq D$
		
		我们假设$F_k$是结合$P$中尺寸最小的文件，则必然有$S_k \geq S_1$。
		
		我们构造另一个集合：$P' = (P - F_k) \cup \{F_1\}$
		
		那么有：
		$$|P'| = |P| - S_k + S_1 \leq |P| \leq D $$
		并且集合$P'$与$P$包含的文件总数相同，

		所以$P'$也是原问题的一个优化解。

		所以我们证明了必然存在某个优化解包含尺寸最小的文件$F_1$
		
		\item \textbf{2. 优化子结构：集合$P$是一个优化解并且包含$F_1$，则集合$Q = P - F_1$必然是$D' = D - S_1$问题的优化解。}
		
		我们假设$Q$不是新问题的优化解，也就是说必然存在一个集合：$R$，满足：
		$$|R| < D'$$
		并且$R$中的文件总数多于$Q$
		
		那么有：
		$$|R \cup {F_1}| < D' + S_1 < D$$
		并且并且$R$中的文件总数多于$P$
		
		这就说明$P$并不是原问题的优化解，与题设矛盾，因此假设错误。

		因此我们证明了集合$Q = P - F_1$必然是$D' = D - S_1$问题的优化解。

		综上所述，原问题可以采取每次“选择最小的文件”的贪心算法解决，因此结论得证。
	\end{enumerate}

	\textbf{2) 解：}

	不是，反例如下：

	$n=4, S_1 = 3, S_2 = 3, S_3 = 3, S_4 = 8, D=9$
	
	这种情况，如果选择文件：$F_1, F_2, F_3$，空间利用率可以达到：$\frac{3+3+3}{9}=100\%$
	
	但是如果按照贪心的策略，第一次选择了文件：$F_4$，剩余空间：$D-S_4=1<3$，不可再放入任何文件，此时空间利用率只有：$\frac{8}{9}=88.9\%$

	因此，这种问题采用贪心策略无法得到全局最优解。
\end{enumerate}



\end{document}
